\seccionreporte{Introducción}{sec:introduccion}

\propositoseccion{%
  Contextualiza el hito del proyecto final: un microservicio de \textbf{aprendizaje no supervisado} que funcionará de forma \textbf{independiente} antes de integrarse a la aplicación móvil.%
}

Entonces, este reporte documenta el primer microservicio de machine learning que desarrollamos para el proyecto de Minería de Datos. En esta etapa todavía no abordamos la integración con la aplicación móvil, sino que nuestro objetivo es demostrar que el backend funciona correctamente de forma independiente al ejecutar el algoritmo no supervisado seleccionado.

Básicamente, el microservicio expone una API REST que espera la información del cliente. En cuanto llegan los datos, primero realizamos una validación estricta para asegurar que el documento contenga lo necesario, luego aplicamos el preprocesamiento, y después de eso procesamos la información con el modelo para retornar el análisis en formato JSON~\cite{han2011data,scikit-learn}. Adicionalmente, guardamos cada una de estas inferencias en una base de datos local para mantener un historial de consulta, y toda la documentación de los endpoints queda expuesta automáticamente mediante Swagger~\cite{openapi}.

\subsection{Alcance de este entregable}

\begin{itemize}
  \item Código fuente del microservicio (\url{https://github.com/eduartrob/integratorProjectClustering-back-corvus}).
  \item Este documento en PDF (compilado desde \LaTeX{}).
  \item Notebook con evidencia de entrenamiento y pruebas.
  \item Diagrama de arquitectura (PDF y PNG).
  \item Documentación de la API (Swagger interactivo).
  \item Colección de pruebas (Postman).
  \item Video demostrativo de 5--10 minutos (\ReporteVideo).
\end{itemize}

\subsection{Modalidades de implementación}

Para esta parte, decidimos utilizar una implementación mixta porque nos pareció el enfoque más lógico para el problema. 

Primero, para el clustering principal utilizamos un \textbf{Modelo local}. Entrenamos HDBSCAN junto con UMAP y los exportamos como archivos \texttt{.joblib}. Entonces el backend los carga directamente en memoria y realiza la inferencia de manera rápida sin depender de servicios externos. 

Posteriormente, para obtener un análisis más profundo, conectamos los resultados con un \textbf{Modelo mediante API} local. Lo que hacemos es enviar los resultados de los clusters a un modelo de lenguaje (Ollama) que se ejecuta en el mismo servidor. Así logramos que el LLM genere recomendaciones precisas basándose en la matemática del modelo local, lo cual optimiza el proceso y evita saturar la RAM con modelos demasiado pesados.

\begin{tcolorbox}[checklist,title={Reutilización en el proyecto final}]
  Nuestro plan es que todo este componente se va a \textbf{conectar directo} con la aplicación móvil en la siguiente fase. Por eso definimos claramente los contratos de la API y preparamos la base de datos de auditoría, asegurando que la integración sea transparente.
\end{tcolorbox}
